几乎所有光学系统都有一个共同任务：收集从物点发出的部分光线，经过反射或折射后，使它们严格或近似会聚于另一个点；或者使出射光看起来像从某一点发出。从波动观点说，光学元件改变入射波面的形状，使各部分在像点具有相同相位。完美成像一般需要复杂非球面和极高加工精度；实际系统往往限制孔径和视场，并在像质、通光量、体积与成本之间折中。

这项任务包含两个不可混淆的要求：一是“同一个物点”的各条光线应在同一个像点重合，二是不同物点的像还应保持可识别的相对位置。前者若不满足会形成模糊斑，后者若不满足会产生畸变。增加孔径能收集更多光、提高衍射分辨率，却会让球面像差更显著；缩小孔径可改善近轴近似，却降低亮度。所谓光学系统的像质，正是在这些条件和加工误差共同作用下评定，而不是只看某一条特征光线是否穿过焦点。

球面容易加工的原因也很直接：两个相配曲面在研磨时，高出的区域优先被磨去，长期相对运动会趋向共同球面。非球面能够更好校正像差，但面形检测和装调更难。图\ref{ch28:fig:contexts}A 的研磨情景解释为什么实际器件偏爱球面；后面的图 20.1 则说明使用球面所付出的代价：一般只能在一定孔径内近似成像。

实际设计还必须限定视场。轴上物点的近轴光也许能很好会聚，离轴物点却可能出现彗差、像散和像面弯曲；因此“一个焦距”只描述系统的一阶性质，并不等于整幅画面处处清晰。镜头组通过不同材料和不同曲率的元件互相补偿这些误差，代价是表面数量、重量、装调难度和成本增加。

成像结论必须同时注明孔径、视场和近轴条件。

\historyfigure{optical-instruments-built-in-v1.png}{成像器件与仪器：光具座、暗箱、透镜、显微镜、望远镜和人眼。}

图\ref{ch28:fig:contexts}把加工、折射现象、成像实验和仪器放在同一条学习线上。A 的玻璃坯与研磨盘以配合球面接触；B 的刚性鸭形模型跨越水面，水下部分的视位置改变；C 的光源、透镜和屏共用一条光具座；D 的反光镜将镜头来光折向上方取景器。下面分别用球面、平面界面和逐级成像解释这些现象，场景图不用于量取角度或焦距。
\historyfigure[0.94\linewidth]{imaging-contexts.png}{研磨球面（A）、水面折射演示（B）、透镜成像实验（C）和单反相机剖面（D）。\label{ch28:fig:contexts}}

\subsection*{球面折射}

透镜是最常见的成像元件，人眼本身也包含折射系统。严格点到点成像通常要求非球面，但球面易于研磨加工；在限制孔径、只取近轴光线后，球面可以近似成像。

图 20.1 中，物点 $S$ 和像点 $P$ 位于球面顶点 $V$ 的两侧，顶点距离分别为 $u,v$；球面半径为 $R$，两侧折射率为 $n_1,n_2$。若入射点 $A$ 对应球心角 $\varphi$，两段几何长度为
\[
 \ell_1=\sqrt{(u+R)^2+R^2-2R(u+R)\cos\varphi},
\]
\[
 \ell_2=\sqrt{(v-R)^2+R^2+2R(v-R)\cos\varphi}.
\]
上述几何先针对 $u,v,R>0$ 的实物实像配置：令 $V=(0,0)$、$C=(R,0)$，则 $A=(R-R\cos\varphi,R\sin\varphi)$，$S=(-u,0)$、$P=(v,0)$。直接计算 $|A-S|$ 和 $|P-A|$ 就得到两式；特别地，$\varphi=0$ 时分别回到 $u,v$。
总光程 $\mathcal L=n_1\ell_1+n_2\ell_2$。对非轴上入射点求导，费马条件给出
\[
 \frac{\partial\mathcal L}{\partial\varphi}
 =R\sin\varphi\left[\frac{n_1(u+R)}{\ell_1}-\frac{n_2(v-R)}{\ell_2}\right]=0.
\]
括号内一般随 $\varphi$ 变化，故不能用一个固定的 $v$ 同时满足所有孔径光线。近轴条件 $|\varphi|\ll1$ 下，令 $\ell_1\to u,\ell_2\to v$，得到
$n_1(1+R/u)=n_2(1-R/v)$，即
\begin{equation}
 \frac{n_1}{u}+\frac{n_2}{v}=\frac{n_2-n_1}{R}.
\end{equation}
将近轴式推广到虚物和虚像时，采用“实正虚负”的物像距约定，而不是把左右坐标都按同一方向读数：发散入射光对应实物 $u>0$，会聚入射光对应虚物 $u<0$；出射光真实会聚对应实像 $v>0$，出射光发散而反向延长线相交对应虚像 $v<0$。$R$ 在球心位于出射侧时为正，位于入射侧时为负。虚共轭须使用有向光程或延长线，不能把上面的正几何长度直接当成负长度。只要在每个界面重新识别物方、像方和传播方向，这一约定可连续使用。

近轴式并不是把严格关系中的角度直接“删掉”，而是保留光程随孔径高度 $y$ 的最低非平凡项。球面在顶点附近的弧高为 $x\simeq y^2/(2R)$。物点到界面的距离与界面到像点的距离分别展开为
\[
 \ell_1\simeq u+x+\frac{y^2}{2u},\qquad
 \ell_2\simeq v-x+\frac{y^2}{2v}.
\]
代入 $\mathcal L=n_1\ell_1+n_2\ell_2$ 后，若所有近轴光线到达像点的光程相同，$y^2$ 项的系数必须为零：
\[
 \frac{n_1}{2u}+\frac{n_2}{2v}+\frac{n_1-n_2}{2R}=0,
\]
这就得到球面折射公式。被略去的 $y^4,y^6,\ldots$ 项正是大孔径时球差等像差的来源，所以“球面能够成像”必须连同近轴、窄孔径条件一起表述。

严格关系中，$\ell_1,\ell_2$ 本身随入射点角度 $\varphi$ 改变；因此即使对某一条光线找到了令光程驻定的入射点，也不能保证所有 $\varphi$ 都给出同一个 $P$。近轴展开的作用，是让最低阶光程差只由 $u,v,R,n_1,n_2$ 决定而不再显含孔径高度。若要让大角度光也严格会聚，就必须改变曲面形状，使从物点经界面到像点的总光程对整个有效口径都相等；这一般给出非球面。例题 20.6 的特殊共轭点则利用几何与折射率关系，在球面上实现了少见的无球差成像。

\SphereRefractionPanel
\VirtualSignPanel

虚像点没有真实能量会聚，不能在该点直接放胶片或 CCD 接收，只能让另一光学系统或眼睛再次处理这些发散光。虚物则不会由普通发光物体单独产生，它只会出现在多元件系统中：前一元件产生的会聚光尚未到达会聚点，就先遇到下一界面。

“虚”描述的是到达当前元件之前或离开当前元件之后的光束状态，而不是说这个点没有物理意义。实物发出的是发散束；虚物对应尚未完成会聚的入射束。实像处确有光能会聚，可以放屏接收；虚像处没有出射光通过，只能由反向延长线定位。判断顺序应当是先看箭头方向，再看光束是会聚还是发散，最后决定 $u,v$ 的正负。仅凭某一点画在元件左边还是右边，在反射系统或折返光路中很容易出错。

例如，把一个新的折射面插入原本会在其后方会聚的光束中，这个会聚点就是该折射面的虚物，$u<0$；折射后若光仍在后方真实相交，则 $v>0$。反过来，普通物体经放大镜后出射光发散，延长线在物方相交，所以 $u>0,v<0$。图 20.2 左图为实物经平界面形成虚像，右图为虚物经平界面形成实像。虚线分别表示出射光的反向延长线和未遇界面时的入射光会聚路径，不是真实传播光线。

\subsection*{眼睛}

把人眼粗略看成空气到折射率 $n=4/3$、曲率半径 $R=5.56\,\mathrm{mm}$ 的单球面系统。远处物体 $u\to\infty$ 时，平行光的像距
\begin{equation}
 f=\frac{nR}{n-1}\approx2.22\,\mathrm{cm},
\end{equation}
与角膜到视网膜的尺度接近。真实人眼还包含角膜、房水、晶状体和玻璃体；水晶体由睫状肌调节曲率，使不同物距的像仍落在固定的视网膜上。水下看不清物体，是因为水的折射率约 $1.33$，与角膜约 $1.376$ 过于接近，角膜折射能力显著下降。

动物眼大体可分为两种工作思路：人眼一类用折射系统在感光面上形成倒立实像；昆虫复眼则由许多小眼单元分别接收不同方向的光。对人眼而言，视网膜位置几乎固定，调焦不能靠前后移动“屏”，只能由睫状肌改变晶状体曲率和等效焦距。看远处时睫状肌较放松；物体移近时，为保持同一像距，系统光焦度必须增大。晶状体形变能力有限，于是出现近点。图 20.3 中角膜、晶状体、视网膜以及光线最终会聚的位置，是理解这一调节链条所必需的四项信息。

\EyePanel

能清楚聚焦的最近点称为近点。青年人的典型明视距离取
\begin{equation}
 L_0=25\,\mathrm{cm}.
\end{equation}
它是目视仪器设计的重要尺度。近点随年龄和屈光状态变化；近视者的远点有限，远视或老视者的近点则较远。$20\,\mathrm{cm}$ 与 $30\,\mathrm{cm}$ 只能作直观示例；实际诊断还需结合远点和屈光度，不能只凭一个距离。

正常眼完全放松时，远处平行光应落在视网膜上。近视眼的系统光焦度相对过强或眼轴偏长，平行光在视网膜前会聚，需要负透镜先使光束稍发散；远视眼或部分老视状态则需正透镜提供额外会聚。显微镜、望远镜的最终像通常设计在无穷远附近，让眼睛放松观察；若使用者近视且不戴镜，调节目镜会把最终虚像移到其有限远点。后面的显微镜例题正是通过“最终虚像应落在哪里”判断镜筒移动方向。

\subsection*{视深度}

平面界面相当于 $R\to\infty$ 的球面折射面，故
\begin{equation}
 \frac{n_1}{u}+\frac{n_2}{v}=0,\qquad v=-\frac{n_2}{n_1}u.
\end{equation}
实物在单一平面折射后成为虚像。观察者在低折射率介质中看高折射率介质里的物体，视深为真实深度乘 $n_2/n_1$，所以水底看起来变浅；反过来从水中看空气里的物体，视距离增大。结论只对近轴光线成立，大视角时不同光线的反向延长线不严格共点。

图 20.4 可直接给出这一比例。物体在介质 $n_1$ 中的真实深度为 $u$，观察者在介质 $n_2$ 中。取一条靠近法线的光，横向位移 $y$ 同时满足
\[
 y\simeq u\theta_1\simeq |v|\theta_2,
 \qquad n_1\theta_1\simeq n_2\theta_2.
\]
消去角度得
\[
 |v|=\frac{n_2}{n_1}u.
\]
从空气看水中，$n_2/n_1\approx1/1.33$，所以视深约为真实深度的四分之三；从水中看空气，比例反过来。这个推导同时说明观察方向若偏离法线很大，$\tan\theta\simeq\theta$ 与 $\sin\theta\simeq\theta$ 不再成立，所谓唯一“视深”只是近似。

\ApparentDepthDiagram

\subsection*{薄透镜}

传统透镜至少包含一个曲面。中心比边缘厚的透镜通常会聚，中心较薄的透镜通常发散。忽略厚度的薄透镜可看作两个球面折射面的组合。对置于空气中的折射率 $n$ 的薄透镜，两个有向曲率半径为 $R_1,R_2$，连续应用球面折射式并消去中间像，得到制镜者公式和成像公式
\begin{equation}
 \frac1f=(n-1)\left(\frac1{R_1}-\frac1{R_2}\right),
 \qquad
 \frac1u+\frac1v=\frac1f.
\end{equation}
会聚透镜 $f>0$，发散透镜 $f<0$。横向放大率（带方向）为 $m=-v/u$，其绝对值
\begin{equation}
 M_T=\left|\frac vu\right|=\left|\frac f{u-f}\right|.
\end{equation}
负号表示实像倒立。作图时可用三条近轴特征光线：过光心近似不偏折；入射平行轴的光经像方焦点；入射通过物方焦点的光出射平行轴。

把两个折射面的中间步骤写出，可以避免把 $R_2$ 的符号背错。第一面由空气进入玻璃：
\[
 \frac1u+\frac n{v_1}=\frac{n-1}{R_1}.
\]
薄透镜厚度忽略后，第一面的像是第二面的虚物，故第二面物距 $u_2=-v_1$。从玻璃射回空气：
\[
 \frac n{u_2}+\frac1v=\frac{1-n}{R_2}.
\]
两式相加时 $v_1$ 恰好消去，才得到制镜者公式。双凸透镜按从左向右传播的约定通常有 $R_1>0,R_2<0$，所以两项共同贡献正光焦度；若只把两个半径都当作无符号长度，第二项就会误消掉第一项。

透镜成像图还提供一套数值自检：$u>2f$ 时为缩小倒立实像；$u=2f$ 时等大；$f<u<2f$ 时为放大倒立实像；$0<u<f$ 时为正立放大虚像。$u\to\infty$ 应给出 $v\to f$，$u\to f^+$ 时 $v\to+\infty$。这些极限若与计算结果不符，应先检查物距、焦距和传播方向的符号，而不是修改作图去迎合错误答案。

“中心厚”与“会聚”是对普通空气中、普通正折射率玻璃的经验说法；决定光焦度的仍是折射率差和两个有向曲率。把同一玻璃透镜浸入折射率更高的液体，光焦度可能减小、为零甚至反号。薄透镜近似还要求厚度相对曲率半径和物像距离足够小，使两个主平面可合并为一个平面；厚透镜中必须保留两面的间隔，中间像位置也不能简单取 $u_2=-v_1$。

图 20.5 下图把两个界面分开：$P$ 处以 $C_1P$ 为法线折射，玻璃内部的 $PQ$ 保持直线；到 $Q$ 才以 $C_2Q$ 为法线再次折射。上面的薄透镜公式，是把这两个依次发生的过程在厚度可忽略时合并，而不是让光线在玻璃内部任意拐弯。
\ThinLensPanel
\LensImagingDiagram

图\ref{ch28:fig:contexts}C 的实验中，移动屏直到光斑最清晰，所测的屏距应与物距一起满足成像公式；若像为虚像，不能靠移动出射侧屏找到它。

例如玻璃折射率 $n=1.54$，双凸透镜两面半径大小分别为 $100\,\mathrm{cm}$ 和 $20\,\mathrm{cm}$（按方向取 $R_1=100\,\mathrm{cm},R_2=-20\,\mathrm{cm}$），则 $f\approx30.9\,\mathrm{cm}$。物距 $u=70.0\,\mathrm{cm}$ 时，$v=fu/(u-f)\approx55.2\,\mathrm{cm}$，$|m|\approx0.788$，形成缩小倒立的实像。

横向放大率来自过光心近似不偏折的两相似三角形，故带符号形式为 $m=-v/u$。它只描述垂直主轴的微小线段，不保证任意平面图形经透镜后保持相似：图形上不同点若有不同轴向坐标，就具有不同物距和放大率。圆、抛物线以及运动轨迹经透镜后是否仍保持原曲线类型，需要把每个点分别按 $X=fx/(x-f)$、$Y=-fy/(x-f)$ 变换，不能把一个统一倍率套到整幅有深度的图形上。这正是例题 20.3 的出发点。

\subsection*{放大镜}

人眼感受到的“大小”由视角决定。尺度 $H$、距离 $L$ 的小物体张角 $\theta\simeq H/L$。光学仪器的角放大率定义为使用仪器后的视角与裸眼基准视角之比。

放大镜是一枚会聚透镜。把物放在焦点以内，透镜形成正立放大的虚像，眼睛靠近透镜观察。若物接近焦点，出射光近似平行，视角 $\theta\simeq H/f$；裸眼把物放在明视距离 $L_0$ 时，$\theta_0\simeq H/L_0$，所以放松眼调焦的角放大率为
\begin{equation}
 M=\frac{\theta}{\theta_0}=\frac{L_0}{f}.
\end{equation}
放大镜的本质，是让眼睛能在小于明视距离的位置清楚观察物体。望远镜、显微镜等目视系统的目镜也起同样作用：它观察前级系统形成的中间像，并进一步放大视角。

若不要求眼睛放松，而把最终虚像调到近点 $v=-L_0$，薄透镜公式给出 $1/u=1/f+1/L_0$。眼睛紧贴透镜时 $\theta\simeq H/u$，相对裸眼近点视角 $H/L_0$ 的角放大率为
\[
 M_{\rm near}=\frac{L_0}{u}=1+\frac{L_0}{f}.
\]
因此教材中常见的 $L_0/f$ 与 $1+L_0/f$ 对应两种不同调焦约定，并不矛盾。前者观察舒适，后者放大略大却更易疲劳。图 20.7 的虚线反向延长线必须交于物体同侧的虚像，不能把放大镜误画成在屏上形成实像。

\MagnifierDiagram

\subsection*{球面镜与平面镜}

非球面反射镜造价较高，球面反射镜在近轴范围内同样可近似成像。图 20.8 中 $CP$ 是法线，反射角等于入射角；对三角形作小角近似可得
\begin{equation}
 \frac1u+\frac1v=\frac2R=\frac1f,\qquad f=\frac R2.
\end{equation}
实物、虚物、实像和虚像的符号意义与透镜一致，只需注意反射后传播方向反转。横向放大率仍为 $m=-v/u$。凸球面镜可取 $R<0$、$f<0$，其作用类似发散透镜。

几何推导中，入射点 $P$、球心 $C$ 和像点必须位于同一子午平面。设物点、像点和球心在主轴上的距离分别为 $u,v,R$，而 $P$ 相对轴的高度很小。半径 $CP$ 是法线，反射定律给出法线两侧角度相等。把三角形中的正弦用小角本身代替，并用“物点方向角 $+$ 像点方向角 $=2\times$ 法线方向角”，就得到
\[
 \frac yu+\frac yv\simeq2\frac yR.
\]
约去非零的 $y$ 后即为球面镜公式。这个过程明确显示了因子 $2$ 来自法线把方向转过两倍，而不是来自光程“往返两次”。

平面镜是 $R\to\infty$ 的极限，因此
\begin{equation}
 v=-u.
\end{equation}
普通实物经平面镜成等大的虚像；一束原本将会聚于镜后虚物点的光，经平面镜也可以在镜前形成实像。单镜头反光相机的折转光路中就包含类似过程。

图\ref{ch28:fig:contexts}D 中反光镜处于镜头预定焦面之前，它把仍会聚的光束折向取景光路，并不是把镜头的会聚能力消除；拍摄时反光镜让开，光才能到达后方感光面。

球面镜公式同样只在近轴范围成立。对凹面镜，平行轴的边缘光与近轴光通常不在同一点会聚，这是球差；把孔径缩小能减轻球差，却同时降低通光量并增强衍射。凹面镜的 $f=R/2$ 只是顶点附近的傍轴焦距，不应把图 20.8 上任意高度的反射点都强行连到这一点。平面镜则没有这一曲率近似：用镜像法可精确证明物、像关于镜面对称，但虚像一侧并无真实光线通过。

镜像法的精确性可由任意一点验证。把物点 $O$ 关于镜面对称到 $O'$，对镜面上任一反射点 $P$ 都有 $PO=PO'$；反射定律又使折线 $O\to P\to$ 观察者等价于从 $O'$ 发出的直线。所有反射光的反向延长线因此严格通过同一个 $O'$，且 $|v|=u$、横向倍率为 $1$。若入射光原本向镜后一点会聚，把该点视为虚物并作同样的对称，反射光会在镜前真实会聚，这说明平面镜也能“虚物成实像”。

\SphericalMirrorDiagram

\subsection*{元件组合}

多元件系统应沿传播方向逐级成像：前一元件的像就是后一元件的物。两薄透镜焦距为 $f_1,f_2$、间距为 $D$，第一透镜物距为 $u_1$ 时
\[
 v_1=\frac{f_1u_1}{u_1-f_1},\qquad
 u_2=D-v_1,\qquad
 v_2=\frac{f_2u_2}{u_2-f_2}.
\]
$v_1>D$ 时，第二透镜前的光正在会聚，故 $u_2<0$ 是虚物；$v_1<0$ 时第一透镜产生虚像，第二透镜看到的物距为 $D+|v_1|$。总横向放大率是各级带符号放大率之积。只背数值而不画传播方向，很容易把虚物和虚像混淆。

\SequentialImagingDiagram

图\ref{ch28:fig:sequential}的 $I_1$ 位于两镜之间，所以第二透镜的物距为正；若把第二透镜移到 $I_1$ 左侧，$I_1$ 就成为第二透镜的虚物。这个正负转换来自光束状态，而不是把中间像删掉。

逐级计算时建议在每个元件前写三行：光线此刻向哪边传播、入射束是发散还是会聚、若没有该元件它将从哪里来或到哪里去。第一问决定正方向，第二问决定实物或虚物，第三问给出物距的几何大小。中间像即使落在两个元件之外，也不能被删除；它仍通过入射束的会聚或发散状态影响下一次成像。最后还应检查总放大率 $m=m_1m_2\cdots$ 的符号与图中倒立次数是否一致。

显微镜和望远镜体现了两种典型组合方式。显微镜的正物镜把近处小物形成放大的中间实像，正目镜再把这个中间像当作焦内实物，给眼睛形成放大的虚像；望远镜的物镜把远处目标形成焦面附近的中间实像，目镜主要放大它对眼睛的视角。两者都需要中间像进入目镜的有效孔径，否则即使成像公式给出位置，观察者也收不到完整光束。Galileo 望远镜用负目镜在物镜光束尚未会聚前截取虚物，系统短且成正像，却没有可放屏或十字丝的中间实像，视场也较受限制。

评价组合系统不能只看倍率。物镜口径决定收集光的能力和衍射分辨率，焦距组合决定角放大率，孔径光阑决定哪些倾斜光束能通过，而像差、镀膜和装调决定这些理论性能能否实现。倍率超过物镜能够提供的空间信息时，只是把衍射斑和像差一同放大。于是“高级望远镜”的简易判断应先看有效口径、像质和机械稳定性，再看可用倍率，而不能只比较广告中的最大倍数。

\Needspace{14\baselineskip}
\section{成像条件与组合系统的计算}

\begin{origexamplebox}{例题 20.1【球透镜】}
空气中半径为 $R$、折射率为 $n>1$ 的玻璃球，在近轴、窄孔径条件下把球外的小明成一个球外实像。求小明到其像的最小距离。
\end{origexamplebox}
\begin{xesolution}
设物到第一球面的距离为 $x$。第一次折射（空气到玻璃，半径 $R$）给出
\[
 \frac1x+\frac n{v_1}=\frac{n-1}{R},\qquad
 v_1=\frac{Rn x}{(n-1)x-R}.
\]
第二球面的物距为 $u_2=2R-v_1$，其有向半径为 $-R$。再次折射后得到球外实像距第二面
\[
 v_2=-R\frac{(n-2)x-2R}{2(n-1)x+(n-2)R}.
\]
小明到像的总距离化简为
\[
 L=x+2R+v_2
 =\frac{2(n-1)(R+x)^2}{(n-2)R+2(n-1)x}.
\]
记分母 $D_x=(n-2)R+2(n-1)x$。实物及球外实像要求 $x>0$、$v_2>0$：当 $1<n<2$ 时，$x>(2-n)R/[2(n-1)]$；当 $n=2$ 时任意 $x>0$；当 $n>2$ 时，$0<x<2R/(n-2)$。
这些区间内均有 $D_x>0$。逐项求导得
\[
 L'(x)=\frac{4(n-1)(R+x)\{D_x-(n-1)(R+x)\}}{D_x^2}
 =\frac{4(n-1)(R+x)[(n-1)x-R]}{D_x^2}.
\]
所以 $x<R/(n-1)$ 时递减，$x>R/(n-1)$ 时递增，而驻点在上述每个允许区间内。因此全局最小值为
\[
 \boxed{L_{\min}=\frac{2nR}{n-1}}.
\]
此时光路关于球心对称；两次折射的符号均已按各自传播方向重新定义。
\end{xesolution}

\begin{origexamplebox}{例题 20.2【奇妙的成像】}
用焦距为 $f$ 的凸透镜和平面镜组合成像。点光源经透镜、镜面反射并再次通过透镜后，像回到光源。光源位于主轴上，平面镜垂直于主轴；若平面镜到透镜的距离为 $L$，求光源到透镜的两个可能距离，并说明真实光源何时可实现。
\end{origexamplebox}
\begin{xesolution}
把平面镜后的返程光路关于镜面展开，问题等价于两枚相同薄透镜相距 $2L$，物和最终像分别位于系统两端、距各自透镜均为 $u$。用空气中光线列向量 $(y,\theta)^{\mathsf T}$，传播矩阵为 $T(d)=\left(\begin{smallmatrix}1&d\\0&1\end{smallmatrix}\right)$，薄透镜矩阵为 $F=\left(\begin{smallmatrix}1&0\\-1/f&1\end{smallmatrix}\right)$。依次相乘得
\[
 FT(2L)F=\begin{pmatrix}1-2L/f&2L\\2(L-f)/f^2&1-2L/f\end{pmatrix}.
\]
再在左右各乘 $T(u)$，总矩阵的右上元为
\[
 B=2u(1-2L/f)+2L+\frac{2(L-f)u^2}{f^2}.
\]
同一物点发出的不同倾角光必须在同一像点重合，所以 $B=0$。乘去分母并分解，回到原点的条件即
\[
 (u-f)\,[Lf-Lu+fu]=0.
\]
故
\[
 \boxed{u_1=f,\qquad u_2=\frac{Lf}{L-f}}.
\]
第一种情况下透镜把轴上点源的光变为平行于轴的光，垂直放置的平面镜令其原路返回。第二种可不用矩阵理解：第一次透镜成像恰在镜面上，即 $v_1=L$，于是 $1/u+1/L=1/f$。镜面反射后，这个镜面上的像点成为返程实物，再经透镜成像回原点。两条不同的光线在镜面相交后互换返程路径，不能说每一条都沿自身原路返回。
\SelfConjugateDiagram
第二根对应有限共轭位置。$u_2$ 要成为镜前有限距离实物位置需 $L>f$；$L<f$ 时第二根对应虚物，不能靠移动真实点光源实现；$L=f$ 时第二根趋于无穷远，有限物距方程只剩 $u=f$。
\end{xesolution}

\begin{origexamplebox}{例题 20.3【透镜成像】}
半径为 $R$ 的圆经焦距为 $f$ 的凸透镜成像。（1）圆心放在哪里时像仍为圆？（2）抛物线能否仍成抛物线？（3）透镜中看到的抛体运动能否仍像抛体运动？
\end{origexamplebox}
\begin{xesolution}
取光轴为 $x$ 轴，物点 $(x,y)$ 的近轴像为
\[
 X=\frac{fx}{x-f},\qquad Y=-\frac{fy}{x-f}.
\]
这里先讨论图形位于包含光轴的子午面、圆心在轴上的情形。$x>0$ 表示物点位于透镜左侧距离 $x$ 处，$X>0$ 表示像点在右侧；两侧横向坐标均向上为正。反解得 $x=fX/(X-f)$、$y=-fY/(X-f)$。圆 $(x-a)^2+y^2=R^2$ 变换后为
\[
 [(f-a)X+af]^2+f^2Y^2=R^2(X-f)^2.
\]
展开后 $X^2$ 的系数为 $(f-a)^2-R^2$，$Y^2$ 的系数为 $f^2$，两者相等且非零才仍为圆，得到
\[
 (a-f)^2=R^2+f^2.
\]
对透镜同侧的实物圆，取
\[
 \boxed{a=f+\sqrt{f^2+R^2}};
\]
另一代数根描述另一侧或虚物配置。

对轴向抛物线 $y^2=4p(x-a)$，同样代入得
\[
 f^2Y^2=4p[(f-a)X+af](X-f).
\]
当 $a=f$ 时右侧化为 $4pf^2(X-f)$，所以 $Y^2=4p(X-f)$：\textbf{存在特殊放置}使抛物线仍成抛物线。一般放置则会变成椭圆或双曲型圆锥曲线。这里讨论各曲线上满足近轴与有效口径的有限部分，不意味着无穷大的图形都能通过透镜。

抛体的\textbf{轨迹形状}与\textbf{运动规律}要分开。若子午面内 $x(t)=x_0+v_xt$，则
\[
 X(t)=f+\frac{f^2}{x_0+v_xt-f},\qquad
 X''(t)=\frac{2f^2v_x^2}{(x_0+v_xt-f)^3}.
\]
当 $v_x\ne0$ 时，纵向像加速度随时间变化，即使像轨迹恰为抛物线，也一般不是匀加速度抛体。若圆心偏离光轴为 $(a,b)$，代入后的圆锥曲线另含 $2bfXY$ 项；要在这个子午面仍为圆，还须 $b=0$。这补足了第一问中圆心的横向位置条件。

题目若不限定图形朝向，还存在另一类简单答案：将圆或抛体轨迹放在\textbf{垂直光轴的同一物平面}内，所有点物距均为 $u_0$，倍率 $m=-f/(u_0-f)$ 为常数。圆自然仍成圆，无须前述特殊轴向位置。若该平面中的抛体坐标为
\[
 y(t)=y_0+v_yt,\qquad z(t)=z_0+v_zt-\tfrac12gt^2,
\]
像坐标为 $Y=my,Z=mz$，故 $Y''=0,Z''=-mg$。忽略光传播时间差，在近轴有效视场内，像确实仍作非退化的匀加速度抛体运动。因此第三问“能否”的回答是\textbf{能}，上述横向平面就是明确例子；其等效加速度通常不等于当地重力加速度。
\end{xesolution}

\begin{origexamplebox}{例题 20.4【小明逞强】}
近视的小明不戴眼镜，与视力正常的小红使用同一显微镜。小红调清楚后，小明应上调还是下调镜筒才能看清？
\end{origexamplebox}
\begin{xesolution}
近视眼不戴镜时远点较近，需要显微镜给出较近的最终虚像。目镜焦距为 $f_e$，中间像在目镜焦内、物距为 $u_e$，虚像距离大小 $|v_e|=f_eu_e/(f_e-u_e)$；要使 $|v_e|$ 变小，应使 $u_e$ 变小。

下调镜筒使物镜靠近标本。对物镜而言，物距更接近焦距，所成中间实像距物镜变大；目镜与物镜间距固定，因此中间像到目镜的距离 $u_e$ 变小，最终虚像变近。故应
\[
 \boxed{\text{下调镜筒（使物镜略靠近标本）}}.
\]
把调节方向写成导数也能自检。镜筒向标本移动 $\varepsilon>0$ 时，物镜物距变化 $\Delta u_o=-\varepsilon$，而
\[
 \frac{\dd v_o}{\dd u_o}=-\frac{f_o^2}{(u_o-f_o)^2}<0,
 \qquad u_e=\ell-v_o,
\]
其中 $\ell$ 为固定的两镜间距。因此 $v_o$ 增大而 $u_e$ 减小。若小红原来按放松眼调到无穷远，而小明远点距目镜为 $F_R$（忽略眼距），所需 $u_e=f_eF_R/(F_R+f_e)<f_e$，与上述方向一致。这里比较两人都采用舒适调焦的状态；若眼睛主动调节或佩戴矫正镜，应重新指定最终像的位置。
实际操作应缓慢并观察侧面距离，避免物镜碰到玻片。
\end{xesolution}

\begin{origexamplebox}{例题 20.5【天文望远镜】}
天文台常见双凸式透射望远镜。（1）为什么不采用凹凸式（Galileo 式）作主要天文望远镜？（2）怎样简易判断望远镜的性能高低？
\end{origexamplebox}
\begin{xesolution}
双凸的 Kepler 结构由正物镜先形成实的中间像，再由正目镜放大。它虽然给出倒像，但视场较大，容易在中间焦面放置十字丝、相机、光谱狭缝等。凹凸 Galileo 结构成正像且镜筒短，却没有可接收的中间实像，出瞳位于系统内部，视场较小，不适合主要天文观测附件。

简易比较不能只看标称“倍数”。在大气视宁度和加工质量相当时，物镜有效口径 $D$ 决定集光能力 $\propto D^2$，并决定衍射角分辨率 $\theta\simeq1.22\lambda/D$；因此先看\textbf{有效口径和像质}，再看焦距、消色差、镀膜、支架稳定性和目镜配置。超过口径所允许的倍率只会把模糊斑放大。
\end{xesolution}

\SpecialAndCatEyePanel

\begin{origexamplebox}{例题 20.6【大角度成像】}
显微镜物镜中有一枚两面均为球面的弯月形厚透镜。左表面 $S_1$ 球心为 $C_1$、半径 $R_1$；右表面 $S_2$ 球心为 $C_2$、半径 $R_2$。两球心都在相应表面的左侧，$C_1$ 又在 $C_2$ 左侧；玻璃对空气折射率为 $n>1$，且 $C_1C_2=R_2/n$。物位于 $C_1$。证明：（1）所有通过有效口径的出射光线的反向延长线严格交于一点 $Q$；（2）$QC_2=nR_2$，并判断像的虚实。
\end{origexamplebox}
\begin{xesolution}
从 $C_1$ 发出的任意光线到达 $S_1$ 时沿该球面的半径入射，所以第一面入射角为零，不发生偏折。于是玻璃内所有光线仍像从 $C_1$ 发出。

把图 20.9 中第二球心取为原点：$C_2=(0,0)$、$C_1=(-R_2/n,0)$，在同侧选 $Q=(-nR_2,0)$。设第二面点 $P=R_2(\cos\varphi,\sin\varphi)$，并记 $a=|P-C_1|$、$b=|P-Q|$。直接计算
\[
 a^2=R_2^2(1+n^{-2}+2n^{-1}\cos\varphi),\qquad
 b^2=R_2^2(1+n^2+2n\cos\varphi)=n^2a^2.
\]
因此 $b=na$。外法向为 $\boldsymbol N=P/R_2$，入射方向 $(P-C_1)/a$，候选出射方向 $(P-Q)/b$。用二维叉积求它们与法线夹角的正弦：
\[
 \sin i=\frac{R_2|\sin\varphi|}{na},\qquad
 \sin r=\frac{nR_2|\sin\varphi|}{b}
 =\frac{R_2|\sin\varphi|}{a}=n\sin i.
\]
两方向的切向分量同号，符合玻璃到空气的折射定律。取通常朝右的有效口径 $|\varphi|<\pi/2$ 时，两方向与外法线的点积均为正，确为向外传播且不发生全反射；还需选取两球面间有正厚度、入射线先穿过 $S_1$ 的口径。
出射方向是从 $Q$ 指向 $P$ 再向右延伸，\textbf{不是从 $P$ 射向 $Q$}。故只有反向延长线在 $Q$ 相交，形成严格无球差虚像，且
\[
 \boxed{QC_2=nR_2}.
\]
证明没有使用小角近似，所以 $C_1,Q$ 是该球面的两个共轭无球差点；后续物镜可利用这束仿佛从 $Q$ 发出的光继续成像。图中到 $Q$ 的虚线表示延长线，在 $Q$ 处不能直接放屏接收此像。
\end{xesolution}

\begin{origexamplebox}{例题 20.7【猫眼探秘】}
门上的猫眼由室外侧焦距 $-1.0\,\mathrm{cm}$ 的凹透镜和室内侧焦距 $3.5\,\mathrm{cm}$ 的凸透镜组成，间距 $2.1\,\mathrm{cm}$。说明为什么室外看不清室内而室内能看清室外；若试图从室外看室内应怎样做？
\end{origexamplebox}
\begin{xesolution}
室外远物先经凹透镜，成为位于凹透镜室外侧 $1.0\,\mathrm{cm}$ 的虚像。对凸透镜，这个物的距离为 $u=2.1+1.0=3.1\,\mathrm{cm}<3.5\,\mathrm{cm}$，故
\[
 v=\frac{3.5\times3.1}{3.1-3.5}\approx-27.1\,\mathrm{cm}.
\]
凸透镜把前一级的小虚像再放大，最终虚像位于可调焦的约 $27.1\,\mathrm{cm}$ 处。这里的“放大”针对中间像，并不表示门外景物的总角放大率大于一；广角门镜通常压缩景物的角尺度，实际视场还受孔径和像差限制。

反向观察远处室内景物时，凸透镜先把平行光会聚于其室外侧 $3.5\,\mathrm{cm}$ 处；凹透镜位于途中，面对的是其外侧 $1.4\,\mathrm{cm}$ 处的虚物，即 $u=-1.4\,\mathrm{cm}$。由 $f=-1.0\,\mathrm{cm}$ 得 $v=-3.5\,\mathrm{cm}$，反向光束在凹透镜后发散，反向延长线交于室内侧距该凹透镜 $3.5\,\mathrm{cm}$ 处。能观察到多大范围，还须检查门孔、两片透镜和眼瞳对斜光束的截取，不能仅凭像距断言视场大小。

反向观察还有更直接的调焦困难：若眼睛几乎贴着室外透镜，虚像到眼睛只有约 $3.5\,\mathrm{cm}$，远小于通常约 $25\,\mathrm{cm}$ 的近点；仅仅贴近孔径、遮挡杂散光不能让正常眼清楚调焦。眼距为 $e$ 时虚像距眼约 $e+3.5\,\mathrm{cm}$，后退可能满足近点条件，却会严重损失视场和接收光量。
从光学原理上，可用适当会聚光学系统重新成像，使最终像处于观察眼可调焦的范围；例如忽略元件间距时，让焦距约 $3.5\,\mathrm{cm}$ 的会聚元件以该虚像为实物，可使出射光近似平行。这只解决调焦，不消除门镜原有的有限孔径、像差与视场限制。由此可见，\textbf{瞳孔接光}和\textbf{眼睛调焦}是两个必须分别满足的条件。
\end{xesolution}

\SilveredAndTrainPanel

\begin{origexamplebox}{例题 20.8【镀银透镜】}
一平凸透镜焦距为 $f$，其平面镀银。凸面一侧距透镜 $2f$ 处放置高为 $H$ 的物，物垂直主轴且下端在轴上。（1）作图判断像的虚实；（2）计算像的位置和大小。
\end{origexamplebox}
\begin{xesolution}
光先通过透镜，遇平面镀银层反射，再反向通过同一透镜。薄元件近似下，透镜光焦度 $1/f$ 被使用两次，平面镜本身不增加光焦度，所以组合等效为焦距
\[
 F=\frac f2
\]
的凹面镜。图 20.11 左图用等效薄平面画两条代表光线，它们返回物方并真实相交，故为实像。逐级追踪也得到同一结论：第一次透镜给 $v_1=2f$；在光尚未会聚前由紧贴的平面镜反射，会聚点关于镜面对称到物方 $2f$ 处。这个点对返程透镜是虚物，$u_2=-2f$，于是 $-1/(2f)+1/v=1/f$，得 $v=2f/3$。第一次倍率 $-1$、平面镜横向倍率 $+1$、返程透镜倍率 $+1/3$，乘积为 $-1/3$。

等效镜的物距 $u=2f=4F$，成像公式给出
\[
 \frac1{2f}+\frac1v=\frac1F=\frac2f,\qquad
 \boxed{v=\frac{2f}{3}}.
\]
像位于凸面外物方 $2f/3$ 处。横向放大率 $m=-v/u=-1/3$，所以
\[
 \boxed{H'=\frac H3},
\]
方向倒立。把镀银面误当成“成一次平面镜虚像”而停止追踪，会漏掉返程的第二次折射。
\end{xesolution}

\begin{origexamplebox}{例题 20.9【累次成像】}
折射率 $n=1.5$、直角边长 $6\,\mathrm{cm}$ 的全反射棱镜上方 $6\,\mathrm{cm}$ 放置物体 $AB$。棱镜右侧 $10\,\mathrm{cm}$ 处为焦距 $f_1=10\,\mathrm{cm}$ 的凸透镜，再向右 $15\,\mathrm{cm}$ 为焦距大小 $10\,\mathrm{cm}$ 的凹透镜。求最终像的位置和横向放大率。
\end{origexamplebox}
\begin{xesolution}
先沿图 20.11 右图把棱镜中的折线光路展开。全反射只把光轴转过 $90^\circ$，不改变会聚性质；棱镜内从上表面经斜面到右表面的轴向路程为 $3+3=6\,\mathrm{cm}$。近轴传播时用约化倾角 $p=n\theta$，高度增量为 $\Delta y=dp/n$，所以该玻璃段的等效空气传播长度是 $6/n=4\,\mathrm{cm}$，不是相位光程 $6n$。故从棱镜右面看，物的等效距离为 $6+4=10\,\mathrm{cm}$；再加棱镜到凸透镜的 $10\,\mathrm{cm}$，有 $u_1=20\,\mathrm{cm}$。

凸透镜：
\[
 v_1=\frac{10\times20}{20-10}=20\,\mathrm{cm},\qquad m_1=-1.
\]
凹透镜位于凸透镜右侧 $15\,\mathrm{cm}$，所以到达它的光原本还要在其右侧 $5\,\mathrm{cm}$ 会聚，即 $u_2=-5\,\mathrm{cm}$。取 $f_2=-10\,\mathrm{cm}$：
\[
 v_2=\frac{(-10)(-5)}{-5-(-10)}=10\,\mathrm{cm},\qquad m_2=-\frac{v_2}{u_2}=2.
\]
因此最终形成在凹透镜右侧
\[
 \boxed{10\,\mathrm{cm}}
\]
处的实像，总横向放大率
\[
 \boxed{m=m_1m_2=-2,\qquad |M|=2}.
\]
负号表示相对棱镜展开后的等效物方向倒立；折回实际空间时还应连同棱镜造成的 $90^\circ$ 方位旋转理解。
\end{xesolution}

